\chapter{Discussion critique}
\section{Arbitrage entre couverture et précision}
Le scénario prudent réduit le risque d'appariement en s'appuyant sur quatre groupes dont les sous-postes concordent mieux. La couverture se situe entre 34,6 et 55,9\,\%. Le groupe alimentaire domine ce sous-panier dans plusieurs profils, si bien que la variation estimée peut être particulièrement sensible à son évolution. Le scénario élargi couvre entre 89,5 et 98,3\,\%, à l'aide de correspondances partielles. Le logement comprend des loyers imputés dans l'ENNVM alors que le suivi IPC des loyers repose sur des locataires. Les frontières du transport, des communications et des loisirs diffèrent aussi entre nomenclatures. Les écarts entre les deux scénarios combinent donc l'effet des postes ajoutés et la redistribution des poids. Ils constituent un test de sensibilité au périmètre.

\section{Collecte urbaine et comparabilité des profils}
En l'absence de séries rurales comparables dans les fichiers analysés, le même vecteur de prix est appliqué à tous les profils. L'écart simulé entre un profil rural et un profil urbain isole le rôle des structures budgétaires conditionnellement à cette hypothèse. Les prix locaux, l'accès aux marchés, la disponibilité des produits et l'autoconsommation peuvent modifier l'exposition effective. Les seuils de quintiles propres à chaque milieu compliquent en outre les comparaisons à niveau de dépense absolu comparable. La couverture des groupes de dépenses ne mesure pas la représentativité géographique.

\section{Portée des résultats et vérifications complémentaires}
Le modèle tient les poids du panier constants et utilise des indices déjà agrégés par division. Il ne permet pas d'estimer une élasticité de demande ou de quantifier les substitutions. Les séries et profils ne fournissent pas l'information nécessaire pour identifier une relation causale entre inflation et comportement ou chiffrer une perte effective de pouvoir d'achat. Les profils publiés ne permettent pas non plus de calculer des intervalles d'incertitude issus du plan de sondage.

Trois validations ont la priorité pour approfondir l'étude. La première consiste à vérifier les appariements de sous-postes, notamment les loyers, la santé, les reclassements entre divisions et les services divers. La deuxième serait de comparer les résultats à des prix effectivement relevés en milieu rural. La troisième consiste à tester le mois de référence, la stabilité des poids, et l'usage éventuel d'une moyenne des prix sur la collecte. L'enquête ne fournit pas dans les tables exploitées de poids mensuels permettant une actualisation observée.
